\documentclass[11pt]{article}
\usepackage[a4paper,margin=27mm]{geometry}
\usepackage{amsmath,amssymb,mathtools}
\usepackage{fontspec,unicode-math}
\setmainfont{Latin Modern Roman}
\setmathfont{Latin Modern Math}
\usepackage{xurl}
\usepackage[hidelinks]{hyperref}
\newcommand{\sref}[2]{\hyperlink{source:#1:#2}{[S#2]}}
\newcommand{\cataloguescope}[1]{\paragraph{Mathematical question.}#1}
\hypersetup{pdftitle={UnsolvedMath Problem 9100001}}
\begin{document}
% UnsolvedMath ID 9100001; source catalogue code AMR-090-0001
\setcounter{section}{9100000}
\section{Odd-prime-power periodicity conjecture}\label{9100001}
{\small\sffamily UnsolvedMath ID 9100001 · Graph Theory}
\subsection{Definitions and mathematical statement}

\paragraph{Shared notation.}$\mathbb N=\{1,2,\ldots\}$, and $\mathbb Z,\mathbb Q,\mathbb R,\mathbb C$ denote the integers, rationals, reals and complex numbers. Any other conventions are specified locally.


Define the Euler up/down numbers $E_n$ for integers $n\ge0$ by
\[
\sum_{n\ge0} E_n\frac{z^n}{n!}=\sec z+\tan z.
\]
For $q\in\mathbb N$ this sequence modulo $q$ is eventually periodic. Its minimal preperiod is
\[
s(q)=\min\{s\in\mathbb Z_{\ge0}:\exists d\in\mathbb N\ \forall n\ge s,\ E_{n+d}\equiv E_n\pmod q\},
\]
and its minimal eventual period is
\[
d(q)=\min\{d\in\mathbb N:\forall n\ge s(q),\ E_{n+d}\equiv E_n\pmod q\}.
\]
The prime-level periods are $d(p)=p-1$ for $p\equiv1\pmod4$ and $d(p)=2(p-1)$ for $p\equiv3\pmod4$.\par\textit{Historical context.} These are the two parts of Ramassamy's original conjecture. A 2026 paper reports a proof of the period formula and the counterexample $s(5^5)=4<5$ to the preperiod formula. \sref{9100001}{3}
\paragraph{Statement recorded in the supplied source.}
For every odd prime $p$, does $s(p^k)=k$ hold for every integer $k\ge1$, and does $d(p^k)=p^{k-1}d(p)$ hold for every integer $k\ge2$? \sref{9100001}{2}
\subsection{Short English statement}
Determine the exact starting point and period of the Euler-number sequence modulo every odd prime power; the original two predictions are recorded with their later source context.
\subsection{Sources}
\begin{description}
\item[\hypertarget{source:9100001:1}{S1}] ULAM.AI and the UnsolvedMath Contributors, UnsolvedMath: A Curated Collection of Open Mathematics Problems, record 9100001 (AMR-090-0001). CC BY 4.0. \url{https://huggingface.co/datasets/ulamai/UnsolvedMath}.
\item[\hypertarget{source:9100001:2}{S2}] Sanjay Ramassamy, \emph{Modular periodicity of the Euler numbers and a sequence by Arnold}, Arnold Mathematical Journal (2018) \S\S1--2, Conjecture 1, pp. 1--3. \url{https://amj.math.stonybrook.edu/html-articles/Files-2015-2024/18-79/index.html}.
\item[\hypertarget{source:9100001:3}{S3}] Berke G\"ule\c{c}, \emph{Modular periodicity of the Euler up/down numbers at odd prime powers}, arXiv:2608.27058v2 (2026) Theorems 5.9 and 6.4. \url{https://arxiv.org/pdf/2608.27058}.
\end{description}
\label{end:9100001}
\end{document}
